\documentclass[11pt]{article}

\usepackage{geometry}
\geometry{margin=1in}
\usepackage{hyperref}
\usepackage{titlesec}
\usepackage{listings}
\usepackage{xcolor}
\usepackage{enumitem}

\titleformat{\section}{\Large\bfseries}{\thesection}{1em}{}
\titleformat{\subsection}{\large\bfseries}{\thesubsection}{1em}{}

\title{IMP CNC Production Tracking Prototype}
\author{Technical Design Review Document}
\date{\today}

\begin{document}

\maketitle
\tableofcontents
\newpage

\section{Project Purpose}
This project is a working prototype of a digital CNC production tracking system for IMP.

The system demonstrates how production tracking could function using a tablet interface, a simulated CNC machine, a supervisor dashboard, and AI-assisted analysis of production events.

The prototype is intended to validate workflow design, architecture, and system behavior before any full production implementation is considered.

The prototype simulates a single CNC production cell using:
\begin{itemize}
\item a simulated HAAS CNC machine
\item a tablet-based operator interface
\item a supervisor dashboard
\item an event-driven backend system
\item AI-assisted analysis of production activity
\end{itemize}

\section{Actual Business Problem}
IMP currently operates multiple HAAS CNC machines used for machining production parts and spare components.

The current workflow is manual:
\begin{enumerate}
\item Operators receive printed technical drawings for parts to be produced.
\item During production, the operator writes the number of completed parts directly on the drawing.
\item Scrap parts, machine downtime, and interruptions are typically not formally recorded.
\item After production is completed, the printed drawing is placed in a folder for later review by management.
\end{enumerate}

Important production information is not systematically recorded:
\begin{itemize}
\item machine runtime
\item machine idle time
\item production interruptions
\item scrap quantities
\item operator interventions
\item machine alarms
\end{itemize}

Management therefore lacks real-time visibility into production activity.

\section{System Scope}
The system is a prototype that simulates a single CNC production cell.

Included:
\begin{itemize}
\item CNC machine simulator
\item backend API server
\item operator tablet interface
\item supervisor dashboard
\item event-based production tracking
\item AI analysis module
\end{itemize}

Out of scope:
\begin{itemize}
\item direct integration with real CNC machines
\item ERP integration
\item multi-machine coordination
\item production scheduling systems
\item automated machine control
\item predictive maintenance algorithms
\end{itemize}

The prototype focuses strictly on:
\begin{itemize}
\item operator workflow
\item machine state tracking
\item event-based production logging
\item real-time monitoring
\item AI-assisted data interpretation
\end{itemize}

\section{Stakeholders and Users}
\subsection{Machine Operator}
The operator is the primary user of the tablet interface.

Operator responsibilities:
\begin{itemize}
\item log into the system
\item select production jobs
\item open technical drawings
\item start machine setup
\item start production cycles
\item pause or resume production
\item report scrap parts
\item enter production notes
\end{itemize}

\subsection{Production Management}
Production management uses the supervisor dashboard to observe:
\begin{itemize}
\item machine state
\item production output
\item interruptions
\item event timelines
\item AI-generated insights
\end{itemize}

\subsection{Technical Reviewer}
Senior software developers may review:
\begin{itemize}
\item system architecture
\item event model design
\item database schema
\item API design
\item implementation strategy
\end{itemize}

\section{System Architecture}
The system follows a modular architecture with four primary components:
\begin{itemize}
\item CNC Machine Simulator
\item Backend API Server
\item Operator Tablet Interface
\item Supervisor Dashboard
\end{itemize}

The backend server acts as the central system controller.

It is responsible for:
\begin{itemize}
\item receiving machine events
\item storing production data
\item managing machine state
\item exposing API endpoints
\item broadcasting real-time updates
\item interacting with AI services
\end{itemize}

The CNC simulator represents a single HAAS CNC machine. Machine ID: \texttt{HAAS\_VF2\_01}

The simulator is controlled through operator actions. It does not run as an uncontrolled autonomous machine.

The frontend application provides two modes:
\begin{itemize}
\item operator interface
\item supervisor dashboard
\end{itemize}

\section{Infrastructure and Runtime Environment}
The prototype runs in a simple local development environment.

Office PC hosts:
\begin{itemize}
\item FastAPI backend server
\item SQLite database
\item machine simulator
\item AI integration module
\end{itemize}

Tablet connects through the local network using a browser.

The frontend is a single web application served over LAN.

This architecture is intentionally simple to support:
\begin{itemize}
\item fast development
\item easy debugging
\item low setup complexity
\item realistic demo behavior
\end{itemize}

\section{Event Engine and Machine Event Model}
The system uses a hybrid event architecture:
\begin{itemize}
\item full event log
\item current machine state table
\end{itemize}

\subsection{Machine states}
\begin{itemize}
\item OFFLINE
\item IDLE
\item SETUP
\item READY
\item RUNNING
\item PAUSED
\item ALARM
\item COMPLETED
\end{itemize}

After a job is completed, the machine briefly enters \texttt{COMPLETED} and then returns to \texttt{IDLE}.

\subsection{Event types}
Operator events:
\begin{itemize}
\item operator\_logged\_in
\item operator\_logged\_out
\item job\_selected
\item drawing\_opened
\item setup\_started
\item setup\_confirmed
\item cycle\_started
\item cycle\_paused
\item cycle\_resumed
\item note\_added
\item scrap\_reported
\item job\_finished
\end{itemize}

Machine events:
\begin{itemize}
\item machine\_state\_changed
\item part\_completed
\item alarm\_triggered
\item alarm\_cleared
\item machine\_reset\_to\_idle
\end{itemize}

\subsection{Event schema}
Every event contains:
\begin{itemize}
\item event\_id
\item timestamp
\item machine\_id
\item event\_type
\item machine\_state
\item job\_id
\item operator\_id
\item reason\_code
\item details
\end{itemize}

Production totals are derived from \texttt{part\_completed} events.

\subsection{State transitions}
Valid flow:
\begin{itemize}
\item OFFLINE $\rightarrow$ IDLE
\item IDLE $\rightarrow$ SETUP
\item SETUP $\rightarrow$ READY
\item READY $\rightarrow$ RUNNING
\item RUNNING $\rightarrow$ PAUSED
\item PAUSED $\rightarrow$ RUNNING
\item RUNNING $\rightarrow$ ALARM
\item ALARM $\rightarrow$ READY
\item RUNNING $\rightarrow$ COMPLETED
\item COMPLETED $\rightarrow$ IDLE
\end{itemize}

Constraints:
\begin{itemize}
\item \texttt{part\_completed} only valid in \texttt{RUNNING}
\item \texttt{cycle\_started} only valid from \texttt{READY}
\item \texttt{cycle\_resumed} only valid from \texttt{PAUSED}
\item \texttt{alarm\_cleared} only valid from \texttt{ALARM}
\item \texttt{job\_finished} only valid when a job is active
\item no direct \texttt{IDLE $\rightarrow$ RUNNING}
\item no \texttt{SETUP $\rightarrow$ COMPLETED}
\end{itemize}

\subsection{Reason codes}
Pause reason codes:
\begin{itemize}
\item TOOL\_CHANGE
\item INSPECTION
\item DRAWING\_REVIEW
\item MATERIAL\_CHECK
\item WAITING\_INSTRUCTIONS
\item OPERATOR\_BREAK
\end{itemize}

Alarm reason codes:
\begin{itemize}
\item TOOL\_WEAR
\item PROGRAM\_STOP
\item FIXTURE\_ISSUE
\item DIMENSION\_CHECK
\item UNKNOWN\_FAULT
\end{itemize}

Scrap reason codes:
\begin{itemize}
\item DIMENSION\_OUT
\item SURFACE\_DEFECT
\item WRONG\_SETUP
\item TOOL\_MARK
\item OPERATOR\_ERROR
\end{itemize}

AI may suggest a reason code from operator notes, but the operator always confirms the final value.

\section{Database Design}
Database engine: \texttt{SQLite}

Core tables:
\begin{itemize}
\item machines
\item operators
\item jobs
\item machine\_events
\item machine\_state
\item scrap\_reports
\item ai\_reports
\end{itemize}

\subsection{machines}
Fields:
\begin{itemize}
\item machine\_id
\item machine\_name
\item machine\_type
\item is\_active
\item created\_at
\end{itemize}

\subsection{operators}
Fields:
\begin{itemize}
\item operator\_id
\item operator\_name
\item pin
\item is\_active
\item created\_at
\end{itemize}

\subsection{jobs}
Fields:
\begin{itemize}
\item job\_id
\item part\_name
\item material
\item target\_quantity
\item drawing\_file
\item status
\item created\_at
\end{itemize}

\subsection{machine\_events}
Fields:
\begin{itemize}
\item event\_id
\item timestamp
\item machine\_id
\item event\_type
\item machine\_state
\item job\_id
\item operator\_id
\item reason\_code
\item details
\end{itemize}

\subsection{machine\_state}
Fields:
\begin{itemize}
\item machine\_id
\item current\_state
\item active\_job\_id
\item active\_operator\_id
\item produced\_count
\item scrap\_count
\item last\_event\_id
\item updated\_at
\end{itemize}

\subsection{scrap\_reports}
Fields:
\begin{itemize}
\item scrap\_id
\item timestamp
\item machine\_id
\item job\_id
\item operator\_id
\item quantity
\item reason\_code
\item note
\end{itemize}

\subsection{ai\_reports}
Fields:
\begin{itemize}
\item report\_id
\item timestamp
\item machine\_id
\item job\_id
\item operator\_id
\item report\_type
\item input\_reference
\item output\_text
\end{itemize}

Design principle:
\begin{itemize}
\item event log preserves history
\item machine\_state supports fast UI updates
\item scrap\_reports keeps scrap structured
\item ai\_reports preserves AI traceability
\end{itemize}

\section{API Design}
The backend exposes:
\begin{itemize}
\item REST API endpoints
\item WebSocket live updates
\end{itemize}

REST responsibilities:
\begin{itemize}
\item operator login/logout
\item job selection/finish
\item drawing access
\item machine setup/cycle/alarm control
\item scrap reporting
\item note submission
\item dashboard data
\item AI requests
\end{itemize}

WebSocket responsibilities:
\begin{itemize}
\item machine state changes
\item produced count changes
\item scrap count changes
\item new event creation
\item AI report creation
\end{itemize}

Operator login must use:
\begin{itemize}
\item operator name
\item PIN
\end{itemize}

Drawings are referenced by file path, not stored as binary in DB.

Alarm trigger must be manually available from the demo.

\section{AI Responsibilities and Limits}
AI acts as an analysis and assistance layer only.

AI responsibilities:
\begin{itemize}
\item operator note structuring
\item reason code suggestion
\item production summary generation
\item natural-language question answering
\item downtime analysis
\end{itemize}

AI inputs:
\begin{itemize}
\item machine event history
\item operator notes
\item scrap reports
\item machine state transitions
\item job information
\end{itemize}

AI must never:
\begin{itemize}
\item change machine states
\item trigger machine cycles
\item start or stop production
\item modify production counts
\item create production jobs
\item override operator decisions
\end{itemize}

All AI outputs are recommendations only and must be saved in \texttt{ai\_reports}.

\section{Frontend System Design}
The frontend is a single React application with two modes:
\begin{itemize}
\item Operator Interface
\item Supervisor Dashboard
\end{itemize}

\subsection{Operator Interface}
Features:
\begin{itemize}
\item operator login with name + PIN
\item job selection
\item drawing viewer
\item machine control panel
\item scrap reporting panel
\item operator notes panel
\end{itemize}

Machine control buttons:
\begin{itemize}
\item Start Setup
\item Confirm Setup
\item Start Cycle
\item Pause Cycle
\item Resume Cycle
\item Finish Job
\end{itemize}

\subsection{Supervisor Dashboard}
Panels:
\begin{itemize}
\item machine status
\item production metrics
\item event timeline
\item runtime metrics
\item AI insight panel
\end{itemize}

Visual style:
\begin{itemize}
\item industrial / professional
\item high contrast
\item large buttons
\item simple layout
\item color-coded states
\end{itemize}

State colors:
\begin{itemize}
\item RUNNING = green
\item PAUSED = yellow
\item ALARM = red
\item IDLE = gray
\item SETUP = blue
\end{itemize}

Frontend uses WebSocket updates for immediate refreshes.

\section{Demo Scenario and Presentation Flow}
The demo uses two short production scenarios.

\subsection{Scenario 1 — Normal Production}
\begin{enumerate}
\item operator login
\item job selection
\item drawing opened
\item setup start
\item setup confirm
\item cycle start
\item simulated part completion
\item finish job
\item machine goes COMPLETED then IDLE
\end{enumerate}

Example job:
\begin{itemize}
\item Support Plate
\item Material: S355
\item Target Quantity: 10
\end{itemize}

\subsection{Scenario 2 — Interrupted Production}
\begin{enumerate}
\item start second job
\item run machine
\item manually trigger alarm
\item operator enters note
\item AI suggests reason code
\item clear alarm
\item resume production
\item report scrap
\item run downtime analysis
\end{enumerate}

Example job:
\begin{itemize}
\item Mounting Bracket
\item Material: Aluminum
\item Target Quantity: 8
\end{itemize}

This demonstrates:
\begin{itemize}
\item digital operator workflow
\item real-time machine visibility
\item structured event logging
\item scrap reporting
\item AI-assisted analysis
\end{itemize}

\section{Project Implementation Plan}
The prototype will be built using AI-assisted development.

Technologies:
\begin{itemize}
\item Python
\item FastAPI
\item SQLite
\item React
\item WebSockets
\item OpenAI API
\item Git
\item VS Code
\item Node.js
\item Python virtual environment
\item npm
\end{itemize}

Development approach:
\begin{enumerate}
\item architecture definition
\item event model design
\item database schema design
\item API specification
\item AI feature definition
\item frontend workflow definition
\item modular code generation and integration
\end{enumerate}

Generate in modules, not as one monolithic app.

Modules:
\begin{itemize}
\item database models
\item backend API foundation
\item event engine
\item machine simulator
\item operator routes
\item production routes
\item AI engine
\item WebSocket update system
\item frontend operator interface
\item frontend dashboard
\end{itemize}

Development phases:
\begin{enumerate}
\item Backend foundation
\item Event engine and simulator
\item Operator interface
\item Supervisor dashboard
\item AI features
\item Demo preparation
\end{enumerate}

Use Git and commit after each completed module.

Success criteria:
\begin{itemize}
\item operator login with name and PIN
\item job selection and drawing display
\item machine state transitions
\item automatic part completion simulation
\item pause and alarm handling
\item scrap reporting
\item real-time dashboard updates
\item AI-generated operational insights
\end{itemize}

\end{document}
